% !TeX root = main.tex
Federated learning (FL) trains a shared model without centralising clients' raw data, but practical edge deployments are shaped by processor speed, memory capacity, network path, and non-IID data. This dissertation argues that heterogeneous FL methods cannot be judged by final accuracy alone: wall-clock time, client trips, staleness, submodel quality, parameter coverage, and slow-device influence all affect whether a method is useful after deployment.

To make these effects measurable, this dissertation develops \textbf{fedctl}, a command-line and browser-facing framework for submitting, deploying, monitoring, and retrieving Flower experiments on a real heterogeneous edge testbed. The testbed contains \(30\times\texttt{rpi4}\) and \(20\times\texttt{rpi5}\) Raspberry Pi workers, provisioned with Ansible and orchestrated through Nomad. \texttt{fedctl} records the submitted run config, deploy config, placement, network impairment, logs, metrics, and artifacts, allowing methods and deployment conditions to be compared without changing application code.

Using this infrastructure, the dissertation evaluates submodel FL, asynchronous FL, and their combination. Under compute heterogeneity, reduced-rate submodels provide a practical middle ground between excluding weak clients and forcing them to train full models: \texttt{FIARSE} gives the strongest performance among heterogeneous methods, while \texttt{HeteroFL} and \texttt{FedRolex} give clearer wall-clock savings. Under straggler-induced heterogeneity, \texttt{FedAsync} is fastest in mixed-device runs, but gives weak devices low aggregate influence; \texttt{FedStaleWeight} improves that influence at the cost of slower target time.

Finally, this dissertation proposes \texttt{FedCover}, a coverage-aware aggregation algorithm for asynchronous submodel FL. Directly composing buffered asynchrony with nested submodels is feasible, but stale partial updates make accepted parameter coverage dynamic. \texttt{FedCover}'s capped inverse-coverage correction improves target and post-warm-up time over coverage-disabled \texttt{Async-HeteroFL}, without simply reweighting slow devices. Overall, the results show that realistic FL evaluation must jointly consider learning dynamics and deployment behaviour, since device capability, network conditions, and update timing fundamentally shape practical performance.